\subsection{Acoustic-Prior-Guided Training Design}
The cross-language emergency-versus-normal evidence is converted into training decisions through a conservative mapping from motivation to implementation. First, we apply \textbf{class-aware augmentation}: augmentation policies are conditioned on the label so that emergency cues are not washed out by globally uniform perturbations. This choice follows the pooled evidence that pitch-energy-related cues carry stable emergency sensitivity.

Second, we add a lightweight \textbf{stats branch} that explicitly models four scalar descriptors:
\[
\{\texttt{pitch\_mean},\ \texttt{pitch\_std},\ \texttt{energy\_env\_std},\ \texttt{pitch\_energy\_corr}\}.
\]
Pitch-related statistics are extracted with YIN-style fundamental-frequency estimation~\cite{decheveigne2002yin}, and the feature pipeline follows standard audio tooling practice with librosa~\cite{mcfee2015librosa}. The branch is auxiliary by design and is not presented as a replacement for the main representation pathway.

Third, training uses a multi-objective form:
\[
\mathcal{L} = \mathcal{L}_{\text{cls}} + \alpha \cdot \mathcal{L}_{\text{aux}}.
\]
Here $\mathcal{L}_{\text{cls}}$ is the primary intent-classification loss, and $\mathcal{L}_{\text{aux}}$ encourages consistency with acoustic priors. The role of $\mathcal{L}_{\text{aux}}$ is to regularize optimization toward emergency-sensitive cues; it is not assumed to guarantee cross-language performance improvements on its own.

If an architecture overview figure already exists in the manuscript, this subsection should reference it for branch placement. Otherwise, the text specification above serves as the implementation-level definition without adding a new complex figure.
